% Prospective support fragment. Not a completed study or reviewed implementation.
% Sources: main-study-settings-supplement.json and protocol source snapshot.
% Render with amsmath, amssymb, booktabs, array and hyperref.
\section{Prospective settings supplement}
\label{sec:prospective-settings}

This supplement records the adopted experiment definitions and their source
provenance before main results. Geometry tables use merged source
\texttt{43db1c62}; additional implementation observations use provisional,
unmerged source \texttt{116d08c7}. Those observations require reconciliation
with the reviewed merge before execution. The accompanying JSON records full
commit identifiers and file SHA-256 hashes. No main-run manifest or outcome
is supplied here.

\subsection{Design coordinates and bounds}

Both spaces use the eight scalar coordinates in
Table~\ref{tab:prospective-scalars}. Space B adds the ten group offsets in
Table~\ref{tab:prospective-groups}; its 55 control points are not independently
optimized. The numerical unit in the tables is millimetres; the executable
coordinate vectors use metres. Every group starts at zero. The unit-box
coordinate is $u_i=(x_i-l_i)/(r_i-l_i)$, with $l_i,r_i$ the stated bounds.

\begin{table}[ht]
\centering\small
\begin{tabular}{lp{0.49\linewidth}rrr}
\toprule
Parameter & Meaning & Lower & Initial & Upper \\
\midrule
MAG & Magnet outer-face recession from the pole-axis rotor surface & 6 & 7 & 15 \\
MH & Magnet radial height & 2 & 7 & 12 \\
MW & Magnet tangential width & 10 & 19 & 25 \\
SW1 & Tooth width & 2 & 4 & 6 \\
SW2 & Slot-opening width & 1 & 2.3 & 4 \\
SW3 & Opening-channel depth from bore-corner chord & 0.5 & 1 & 1.5 \\
SW4 & Slot-axis back-iron depth & 2 & 8.25 & 20 \\
Sr1 & Slot-bottom fillet radius & 0.5 & 1 & 2 \\
\bottomrule
\end{tabular}
\caption{Free scalar coordinates in millimetres, shared by spaces A and B.
The initial values are the optimization reference, not the geometry
functions' defaults (which differ for SW3 and Sr1).}
\label{tab:prospective-scalars}
\end{table}

The frozen coordinates are axial length $100\,\mathrm{mm}$, outer stator
diameter $135\,\mathrm{mm}$, radial scale $1$, and operating angle $0^\circ$.
The upstream optimization's wider bounds on those four coordinates are not
part of this study. At zero offsets, the reference rotor and bore radii are
$44$ and $45\,\mathrm{mm}$, with the coupling circle at $44.5\,\mathrm{mm}$.
For the rectangular magnet in the rotor frame,
$x\in[44-\mathrm{MAG}-\mathrm{MH},44-\mathrm{MAG}]$ and
$|y|\leq\mathrm{MW}/2$, in millimetres.

\begin{table}[ht]
\centering\small
\begin{tabular}{lrp{0.43\linewidth}rr}
\toprule
Group & Points & Controlled feature & Lower & Upper \\
\midrule
CPR1 & 2 & Sector-edge surface endpoints & -1.5 & 0 \\
CPR2 & 2 & Edge-arc middle controls & -1.5 & 0 \\
CPR3 & 6 & Bridge-arc controls & -1 & 0 \\
CPR4 & 2 & Pole interior pair 1 & -1.5 & 0 \\
CPR5 & 2 & Pole interior pair 2 & -1.5 & 0 \\
CPR6 & 2 & Pole interior pair 3 & -1.5 & 0 \\
CPR7 & 2 & Pole interior pair 4 & -1.5 & 0 \\
CPS1 & 18 & Opening corners and opening-arc middles & 0 & 0.5 \\
CPS2 & 12 & Tooth-arc middle controls & 0 & 0.5 \\
CPS3 & 7 & Tooth-centre bore controls & 0 & 0.5 \\
\bottomrule
\end{tabular}
\caption{Additional grouped offsets in space B, in millimetres. Rotor groups
contain 18 distinct controls and stator groups 37. All initial offsets are zero.}
\label{tab:prospective-groups}
\end{table}

A group's offset moves each of its controls along that control's frozen
setup direction. These directions do not track later scalar changes; stator
setup uses SW2=$2.3\,\mathrm{mm}$. Weights retain the generator's value at the
current scalars. Rotor controls are ordered by polar angle: CPR1 is
$\{0,17\}$, CPR2 is $\{1,16\}$, CPR3 is $\{2,3,4,13,14,15\}$, and CPR4--7
are the mirror pairs $\{5,12\}$ through $\{8,9\}$.
Stator groups repeat over the six slots, sharing interface controls.
These groups are GEODE's reconstruction from the pinned exported geometry,
group sizes, symmetry and displacement data; they do not claim reproduction
of the unavailable upstream multipatch numbering.

The existing scalar and control-point maps have different discrete motion
conventions. The scalar extension uses base-area weights $A_e^{-1}$;
the control-point extension uses the plain harmonic weights ($\chi=0$).
They also preserve different boundary parameters, so even zero group offsets
do not imply identical moved node coordinates at changed scalars.
Both use their documented sector-edge constraints. Thus a space comparison
includes these discrete motion choices as well as additional shape freedom;
final refinement checks are necessary, and the comparison is not a pure
measurement of intrinsic dimension. The main implementation must record the
actual motion options rather than infer them from the space label.

\subsection{Models, normalization and solver settings}

A/B/C train on L1 and D/E on L0, with coupling torque, $h=45$, twenty rotor
angles, the same nonlinear material law and current fractions $0,0.5,1$ of
$9.6777\,\mathrm{A}$. Each resolution supplies its own $F_r(\theta_0)$ and
$\mu_{r,1}(\theta_0)$. The normalized objective is
$f_r=F_r/F_r(\theta_0)$ and mean residual
$g_r=\mu_{r,1}/\mu_{r,1}(\theta_0)-1.0025$.
The common validator is L2/$h93$/g20, using the same exact anchor for all
methods. The main study must retain the numerical material/solver options,
mesh hashes and model constants in its frozen manifest; this model label
alone is not a reproducible configuration.

\begin{table}[ht]
\centering\small
\begin{tabular}{lll}
\toprule
Setting & SLSQP (A/E) & COBYQA (B/D) \\
\midrule
Field derivatives & Adjoint objective and mean & None \\
Coordinate bounds & Unit box & Unit box \\
Mean and geometry constraints & Explicit & Explicit with NaN refusal barrier \\
\texttt{ftol} option & $10^{-15}$ & Not used \\
Initial/final trust radius & Not applicable & $0.1/10^{-6}$ \\
Initial interpolation set & Not applicable & $2n+1$ evaluations \\
Refused point & Abort, restore, fresh restart & NaN extreme barrier \\
\bottomrule
\end{tabular}
\caption{Registered optimizer choices, with $n=8$ or $18$. Within each
optimizer, settings are identical across spatial resolutions. SciPy 1.18.1
is pinned; success flags never determine target attainment.}
\label{tab:prospective-optimizer-settings}
\end{table}

The provisional driver additionally sets SLSQP \texttt{maxiter=100000},
COBYQA \texttt{maxiter} and \texttt{maxfev} to $10^9$, and \texttt{scale=False}; the explicit
unit-box transform already scales its variables. These are finite
implementation caps, not literally unlimited iterations. These settings were accepted before main outcomes; the reviewed
merge must be checked against them. Every main-run record must state whether
any cap was reached. This does not assert that a cap has been reached.

SLSQP's refusal mechanism permits at most twenty restarts and twenty
restoration bisections, from the last admissible evaluated anchor toward
the rejected point. Geometry-only restoration tests the actual restored
point; mesh-infeasible restoration also tests mesh policy. No admissible
probe ends the run as restoration failure. Fresh repair tries default,
\texttt{Mesh.Algorithm=5} (or 1 when the default is 5), size factor $0.9$,
then $0.8$, each from defaults
and with every cost charged. All-four failure is mesh infeasibility, not
an infrastructure error. The actual algorithm must be recorded and a fixture
must show distinct mesh bytes on attempt 2. Credit and best-feasible selection
allow $h_i\geq-10^{-12}$; optimizer constraints remain $h_i\geq0$ and
buildability checks remain exact. The $5^\circ$ mesh floor is also an implicit
geometric constraint. Space B's unguarded slot-opening wall can narrow to
$0.16\,\mu\mathrm{m}$ at minimum SW3 and maximum CPS1; repair/refusal handles
this known risk without an additional thickness constraint.
Complete strict/discrete guard formulas, scales,
safe downstream derivatives and their regression evidence remain to be
bound to the reviewed implementation.

\subsection{Empirical surrogate algorithm}

Method C fits four cubic RBFs with affine polynomial tails: the three
training torque standard deviations divided by $\sqrt{F_r(\theta_0)}$, and
the full-current mean divided by $\mu_{r,1}(\theta_0)$. For each scalar output,
\begin{equation}
 s(u)=\sum_{j=1}^{m}\lambda_j\lVert u-u_j\rVert_2^3+a_0+a^Tu,
 \qquad \sum_j\lambda_j=0,\quad \sum_j\lambda_j u_j=0.
\end{equation}
All evaluated training points enter the model. Its objective is the sum of
squared predicted normalized standard deviations, and its mean constraint
is $\widehat\mu\geq1.0025$.

The initial set is the start plus one buildable coordinate step per free
variable, at radius $\Delta_0=0.1$. A refused step tries the opposite sign,
then up to three halvings in each direction, skipping out-of-box points.
Failure to obtain a full-rank set ends the run as restoration failure;
refused points are never fitted. After the initial fit, the best feasible-first
incumbent is selected and validated if it differs from the start.

The adopted iteration is:
\begin{enumerate}
\item Minimize the surrogate objective with SLSQP from the incumbent in
$[0,1]^n\cap\{u:\lVert u-u_k\rVert_\infty\leq\Delta\}$, subject to the predicted
mean constraint and real geometry constraints.
\item If the proposed point is within $10^{-3}\Delta$ of an evaluated point,
substitute a coordinate geometry point, cycling over coordinates.
\item Evaluate a buildable candidate on the training FEM; save the old
model's prediction before refitting. Refused candidates are skipped, never
fitted, and count as rejections.
\item Accept a candidate that improves feasible-first ordering. Let
$\rho$ be actual normalized objective reduction divided by predicted
reduction. If $\rho\geq0.75$ (or the old incumbent was infeasible) and the
step is at least $0.9\Delta$, increase the radius to
$\min(2\Delta,0.5)$. Validate accepted incumbents under the common policy.
\item Otherwise halve the radius. The registered geometry-point rule adds
one coordinate point after every second consecutive rejection. This includes
refused candidates: their early return must not skip this step. The added
geometry point is itself checked for buildability and skipped if refused. A
non-improving geometry substitute is itself a rejection; all refusal and
geometry-point edge cases must agree with the reviewed implementation.
\end{enumerate}

Stop at the common wall-time budget, validated T99, restoration failure,
a real evaluator failure, or $\Delta<10^{-6}$. A distinct final non-start
incumbent is validated if it has not already been validated. Training
forwards are charged as forwards, the initial fit as setup/calibration,
and subsequent subproblems/refits as surrogate updates. Nothing is removed
from the charged wall time as overhead.

For reproducibility, the provisional implementation orders points using
\begin{equation}
 v(u)=\max(0,-g_r(u))+\sum_i\max(0,-h_i(u)),
\end{equation}
with failed discrete guards assigned infinite violation. Smaller violation
wins; between feasible points, smaller training $F$ wins. It solves the RBF
saddle system directly, falls back to a flagged least-squares solution if
singular, uses $\rho=-\infty$ for nonpositive predicted reduction, and
solves its surrogate subproblem with
\texttt{maxiter=200}, $\mathrm{ftol}=10^{-14}$. The finite subproblem caps were accepted before main outcomes. The other
observed details still require reviewed freezing; they are not new permissions to alter the
registered comparison after outcomes.

\subsection{Execution provenance preparation}

Private implementation \texttt{d1bae200} adds recorded evaluator builds,
linked recovery with frozen calibration and runtime inputs, and explicit solver
termination records. Its optimizer BLAS/LAPACK fingerprint matches the peer
helper at \texttt{1a3a94ea}; full backend dependency dictionaries also retain
ABI details. The source and evidence bindings are listed separately in
\path{main-study-runtime-preparation.json}. They do not replace the earlier
numeric-table source snapshot or establish an executed main study.

The final execution manifest must identify the actual build, host, concurrency,
NumPy/SciPy and mesher runtimes, BLAS/LAPACK configurations, thread/environment
settings and calibration. Build receipts and controlled-process tests are
preparation evidence; the reviewed native build and main measurements remain
pending. Whole-run solver caps and local surrogate-subproblem caps will be
reported separately, with no change to target credit.

\subsection{Outstanding verification}

The independent source inventory and prospective protocol do not replace
an executable study. Before this supplement becomes publication evidence,
bind the reviewed implementation, full guard/scale and numerical model
tables, exact main starts, host/environment/calibration, gradient sweeps,
mixed-resolution smoke and recovery checks. All main runs and final
candidate checks remain pending. [PENDING reviewed implementation/settings
freeze, complete executable manifests and measured results]
